\section{Introduction}
\label{sec:introduction}

Let $A$ be a finite-dimensional algebra and let
\[
 \MM=\Mor(\proj A)
\]
be the exact category whose objects are morphisms between finitely generated
projective right $A$-modules and whose conflations are degreewise split short
exact sequences.  Bautista's analysis of this category shows that it is
Krull--Schmidt and hereditary as an exact category, even when $A$ itself has
infinite global dimension \cite{Bautista}.  Li proved that its basic maximal
rigid objects are
\[
 C_M\oplus Z_P\oplus K_A
\]
for basic support $\tau$-tilting pairs $(M,P)$, that every rigid object can be
completed, and that the maximal-rigid exchange graph is the support-$\tau$-
tilting graph \cite{Li}.  These results make $\MM$ a natural exact model for
two-term silting combinatorics beyond the hereditary case.

The Hall and cluster structures carried by this model are nevertheless
subtle.  For a finite acyclic quiver, Ding--Xu--Zhang and Fu--Peng--Zhang
construct direct morphism-Hall realizations of quantum cluster algebras
\cite{DingXuZhang,FuPengZhang}.  Zhang's fixed-size complex construction gives
the general exact integration needed for $\Mor(\proj A)$ over a finite field
\cite{ZhangFixedSize}.  However, the canonical exact integration records the
source and target projective terms, not the differential.  For a
nonhereditary cluster-tilted algebra, different maps with the same terms can
have different homology, self-extension spaces and cluster variables.  This
is the differential-blindness phenomenon isolated in the preceding
iteration of this project.

The same preceding work established a compensating rigidity theorem.  In a
fixed affine presentation space
\[
 \Pres(P^-,P^0)=\Hom_A(P^-,P^0),
\]
the normal space to the base-change orbit of a map $f$ is
$\Ext^1_{\MM}(F_f,F_f)$.  Hence rigid maps are exactly the open orbits, and
there is at most one rigid isomorphism class with fixed terms.  Full morphism
homology recovers the reduced object and its split index.  On the union of
rigid-open-orbit cones, this gives a piecewise-unimodular transport
\[
 \Gamma_T:\lvert\Fan_\tau(A)\rvert_{\mathrm{reach}}
 \longrightarrow
 \lvert\Fan_{\mathrm{cl}}(T)\rvert
\]
from the two-term presentation fan to the cluster $g$-fan.  Davison--Mandel's
quantum theta theorem then assigns the path-independent quantum element
\[
 \ThetaMor_{T,q}(\delta)
 =\vartheta_{\Gamma_T(\delta)}
\]
to every rigid presentation weight \cite{DavisonMandel}.  The geometric
importance of walls and chambers in $\tau$-tilting theory has also been
emphasized recently in the construction of the $\tau$-cluster morphism
category \cite{SchrollTattarTreffingerWilliams}.

The last iteration promoted this object-level reconstruction to
multiplication when the mutable exchange matrix $B_0$ is invertible.  Let
$K=K_0^{\mathrm{sp}}(\proj A)$ and
$L=K\oplus K=K_0(\MM)$, and write
\[
 D=(-I\ I):L\longrightarrow K,
 \qquad D(p,q)=q-p.
\]
If $B_0^T\Lambda_0=dI$, the term-lattice form
\[
 \widehat\Lambda_T=D^T\Lambda_0D
\]
corrects the exact skew Euler form $\Lambda_{\MM}$.  Twisting the Hall product
by the resulting bicharacter is associative on the whole exact Hall algebra,
and on every rigid chamber it is isomorphic to the quantum affine-space
algebra of cluster monomials.  The proof used full-rank uniqueness: two
compatible skew forms for the same invertible exchange matrix must coincide.
The singular-rank case was left as a collection of chamberwise forms.

The present paper shows that full-rank uniqueness was convenient but not the
correct organizing principle.  The correct invariant is the holonomy of the
piecewise-unimodular presentation-to-cluster transport.

\subsection*{The transfer group and globalizable quantisations}

Let
\[
 R=f_1\oplus\cdots\oplus f_n
\]
be a reachable ordered basic reduced maximal rigid morphism object.  Its two
integral bases are
\[
 \Delta_R=(I(f_1)\ \cdots\ I(f_n)),
 \qquad
 G_R=(g_T(f_1)\ \cdots\ g_T(f_n)).
\]
Both matrices are unimodular.  We define the chamber transfer
\[
 \mathsf L_R=G_R\Delta_R^{-1}\in\mathrm{GL}_n(\mathbb Z).
\]
It is the local matrix of $\Gamma_T$.

The complete affine family of principal-compatible quantum forms is
parameterized by an initial integral skew form $S$.  Its mutable block in the
cluster chart $R$ is
\[
 \Lambda_R^{\mathrm{mut}}(S)=G_R^TSG_R.
\]
On the Hall side, a global target form which kills contractible classes and
has initial boundary value $S$ is forced to be
\[
 D^TSD.
\]
Its restriction to the reduced term lifts of the chamber is
$\Delta_R^TS\Delta_R$.  Equality of the Hall and cluster exponents is
therefore equivalent to
\[
 \mathsf L_R^TS\mathsf L_R=S.
\]
This gives the fixed lattice
\[
 \Qglob
 =\{S\in\Skew_n(\mathbb Z):
       \mathsf L_R^TS\mathsf L_R=S\text{ for all reachable }R\}.
\]
The result both generalizes and sharpens the full-rank theorem.  It gives a
necessary and sufficient condition, not merely a construction.

The zero form belongs to $\Qglob$ for every exchange matrix.  It corresponds
to the canonical principal-coefficient quantisation, whose mutable--mutable
block is zero at every seed.  Thus one global exact-Hall cocycle exists for
the canonical principal atlas in every rank.  In the full-rank
coefficient-free situation, the compatible form $\Lambda_0$ also belongs to
$\Qglob$, recovering the earlier result.  More generally, the dimension and
integral basis of $\Qglob$ are computable by a finite exact linear system
whenever a finite rigid atlas is available.  This is a fan-holonomy form of
the quantisation-space problem studied by Gellert--Lampe
\cite{GellertLampe}.

\subsection*{Symplectic wall transvections}

Adjacent chamber maps agree on their common codimension-one face.  Their
difference has rank at most one.  After passing to a relative transition, we
prove the integral normal form
\[
 T=I+v\ell^T,
 \qquad \ell^Tv=0,
\]
where $\ell$ is the primitive covector vanishing on the wall.  For any skew
form $S$,
\[
 T^TST=S
 \quad\Longleftrightarrow\quad
 Sv=c\ell
 \quad(c\in\mathbb Z).
\]
This is an edge-local globalization test.  When $S$ is nondegenerate, every
nonzero wall preserving $S$ is an integral symplectic transvection.  In
particular, the full-rank presentation-to-cluster map is not merely
piecewise-unimodular: it is piecewise symplectic.

\subsection*{Mutable globalization versus frozen globalization}

The canonical principal coefficient matrix exists in every rank, but it
contains frozen directions.  A single global form on the enlarged Hall term
lattice with a fixed frozen basis is a stronger requirement than globalizing
its mutable block.  We derive the exact additional equation
\[
 (\Delta_R^T+G_R^T)(dI+SB_0)=0.
\]
This yields a sharp dichotomy.  For the canonical principal parameter $S=0$,
the mutable cocycle always globalizes, whereas the full fixed-frozen form
globalizes only if $\Delta_R=-G_R$ in every chamber.  In a genuine nontrivial
atlas this normally fails.  In the scalar-compatible full-rank case,
$S=\Lambda_0$ satisfies $dI+\Lambda_0B_0=0$, so the mixed obstruction
vanishes and the whole fixed-frozen form globalizes.  Its frozen directions
are central.  Thus ``principal coefficients remove the rank obstruction'' and
``one fixed frozen Hall lattice realizes every seed'' are different
statements; the first is true, the second is governed by the equation above.

\subsection*{Incompatible products as a kernel problem}

The cocycle-corrected Hall-to-theta map is an algebra isomorphism on each
rigid chamber and agrees on common faces.  It is not known to extend to
products of incompatible rigid objects, whose Hall products contain nonsplit
and often non-rigid middle terms.  We replace this vague obstruction by a
universal criterion.  Let $\mathscr U_T$ be the colimit of the rigid face
algebras.  It maps to the subalgebra generated by rigid classes in the twisted
exact Hall algebra and also to the reachable quantum theta algebra.  A global
homomorphism extending every chamber map exists if and only if
\[
 \ker\pi_{\mathrm H}\subseteq\ker\pi_\vartheta.
\]
The inclusion is the precise remaining problem.  Exchange and wall relations
are visible elements of the theta kernel; the missing theorem is that every
relation forced by exact Hall multiplication is already killed there.

We also prove a recognition theorem for the direct morphism-to-critical
comparison.  Keller--Yang mutation, Hall scattering diagrams and
cohomological Donaldson--Thomas theory already provide mutation-covariant
$3$-Calabi--Yau structures \cite{KellerYang,BridgelandScattering,Davison}.
What remains is a canonical class attached to a concrete projective
morphism.  If such an assignment has the correct initial normalization and is
mutation covariant, then its integration is forced on every reachable rigid
object to equal $\vartheta_{\Gamma_T(I(f))}$.  The theorem does not construct
the class, but it reduces the comparison to a local covariance check.

\subsection*{Computation and scope}

For the nonhereditary cluster-tilted algebra
\[
 A=\kk Q/(bc,cd,db),
 \qquad
 Q:1\longrightarrow2\longrightarrow4\longrightarrow3\longrightarrow2,
\]
exact code enumerates the entire paired $A_4$ pattern.  It finds
\AIVTransferMatrices\ distinct transfer matrices, computes the invariant
skew lattice, verifies all \AIVOrientedMutations\ labelled wall transitions,
and solves the full fixed-frozen affine system.  Separate regressions test the
general formulas in full ranks $2,4,6$ and singular ranks $3,5$.  The code is
a certificate for the implementation and the finite example; it does not
replace the general proofs.

No claim is made that incompatible rigid Hall products are closed, that the
whole exact Hall algebra maps to the quantum cluster algebra, or that a
framed critical Hall class has been constructed.  The novelty claims are
therefore deliberately limited to the globalization criterion, its
symplectic wall form, the fixed-frozen obstruction, the universal kernel
criterion, and the resulting exact computations.  The manuscript is a
research draft and does not assert priority beyond the literature search
recorded in \cref{sec:audit}.
